\bookchapter{บทที่ 63: วิว สเตต และการนำทาง}

\chapterintro{ของเล่นแอปหน้าเดียวเราทำเป็นแล้ว บทนี้จะแก้สองปัญหาของการ'โตขึ้น': วิธีกระโดดไปมาระหว่างหลายหน้า (การนำทาง) และข้อมูลไหลระหว่างหน้าอย่างไร (การจัดการสเตต) เรียนจบบทนี้ คุณจะสร้างโครงแอปหลายหน้าสุดคลาสสิกอย่าง'รายการ รายละเอียด แก้ไข'ได้}

\booksection{'ภาพรวมการไหลของข้อมูล'ในการจัดการสเตต}

เครื่องมือจัดการสเตตของ SwiftUI ไม่ได้มีแค่ @State ตัวเดียว ลองดูภาพรวมนี้ก่อน เพื่อเข้าใจว่า'เมื่อไหร่ใช้ตัวไหน':

\begin{center}
\begin{tabularx}{\linewidth}{|X|X|X|X|}
\hline
\textbf{เครื่องมือ} & \textbf{ดูแลอะไร} & \textbf{ใครใช้} & \textbf{สถานการณ์ทั่วไป} \tabularnewline
\hline
@State & สเตตชั่วคราวส่วนตัวของแต่ละหน้า & วิวปัจจุบัน & ตัวนับ, ข้อความในช่องกรอก \tabularnewline
@Binding & ยืมสเตตของคนอื่นมาใช้ & ให้วิวลูกแก้ค่าของวิวแม่ & คอมโพเนนต์ที่สร้างเอง \tabularnewline
@StateObject / @Observable & โลจิกที่ใช้ร่วมกันข้ามหน้า & วิวโมเดล (ViewModel) & การโหลดข้อมูล, โลจิกของแอป \tabularnewline
@EnvironmentObject / .environment & สเตตที่แชร์ทั้งแอป & ทั้งแอป & สถานะล็อกอินของผู้ใช้, การตั้งค่าธีม \tabularnewline
\hline
\end{tabularx}
\end{center}

\begin{center}\small\color{gray}ภาพรวมเครื่องมือจัดการสเตตของ SwiftUI (ปี 2026 @Observable ขึ้นเป็นกระแสหลักแล้ว)\end{center}

\begin{lstlisting}[language=Swift]
// สไตล์ทันสมัยบน iOS 17+: แมโคร @Observable (แนะนำสำหรับโปรเจกต์ใหม่)
import Observation

@Observable
class CounterModel {                 // วิวโมเดล
    var count = 0
    func increment() { count += 1 }
}

struct CounterView: View {
    @State private var model = CounterModel()   // ถือโมเดลไว้

    var body: some View {
        VStack {
            Text("คุณแตะไป \(model.count) ครั้ง")
            Button("แตะฉัน") { model.increment() }
        }
    }
}
\end{lstlisting}

คำแนะนำแบบตรง ๆ สำหรับมือใหม่: อันที่เป็นส่วนตัวของหน้าใช้ @State ส่วนที่ข้ามหน้าหรือมีโลจิกให้ใช้คลาสแบบ @Observable @StateObject/@ObservableObject ในบทเรียนเก่าเป็นสไตล์รุ่นก่อน อ่านออกก็พอ โค้ดใหม่ให้ทำตามข้างบน

\booksection{NavigationStack: 'กองจาน'ของการเปลี่ยนหน้า}

ตั้งแต่ iOS 16 เป็นต้นไป การนำทางหลายหน้าใช้ NavigationStack — เหมือนกองจาน ดันใส่ทีละใบ ย้อนกลับก็เด้งออกทีละใบ:

\begin{lstlisting}[language=Swift]
struct HomeView: View {
    var body: some View {
        NavigationStack {
            List {
                NavigationLink("สมุดของซานผวน") {
                    DetailView(name: "สมุดของซานผวน")   // หน้าปลายทาง
                }
                NavigationLink("สมุดของลี่ซือ") {
                    DetailView(name: "สมุดของลี่ซือ")
                }
            }
            .navigationTitle("รายชื่อสมุดบัญชี")   // หัวเรื่องใหญ่ด้านบน
        }
    }
}

struct DetailView: View {
    let name: String
    var body: some View {
        Text("หน้ารายละเอียดของ \(name)")
            .navigationTitle(name)
            .navigationBarTitleDisplayMode(.inline)
    }
}
\end{lstlisting}

NavigationLink ก็คือ'ลิงก์เปลี่ยนหน้าที่แตะได้' ส่วน navigationTitle สร้างชื่อบนแถบนำทางและปุ่มย้อนกลับให้อัตโนมัติ ถ้าอยากเปลี่ยนหน้าด้วยโค้ด (เช่นกด'บันทึก'แล้วกลับหน้าเดิมอัตโนมัติ) ให้ใช้ @Environment(\textbackslash{}.dismiss) ปิดหน้าปัจจุบัน — สองเทคนิคนี้คู่กัน ก็ครอบคลุมความต้องการด้านนำทางถึง 90\%

\booksection{ป๊อปอัปและแท็บด้านล่าง: sheet กับ TabView}

นอกจาก'ดันหน้าใหม่เข้ามา' ยังมีรูปแบบการโต้ตอบคลาสสิกอีกสองแบบ:

\begin{lstlisting}[language=Swift]
// 1. ป๊อปครึ่งจอ/เต็มจอจากด้านล่าง (ฟอร์มบันทึกชอบใช้)
struct HomeView: View {
    @State private var showAdd = false
    var body: some View {
        Button("จดรายการ") { showAdd = true }
            .sheet(isPresented: $showAdd) {
                AddRecordView()     // หน้าบันทึกที่ป๊อปขึ้นมา
            }
    }
}

// 2. แถบแท็บด้านล่าง ("โครง" ของแอป)
struct MainTabView: View {
    var body: some View {
        TabView {
            RecordListView()
                .tabItem { Label("สมุดบัญชี", systemImage: "book") }
            StatsView()
                .tabItem { Label("สถิติ", systemImage: "chart.pie") }
            SettingsView()
                .tabItem { Label("ตั้งค่า", systemImage: "gear") }
        }
    }
}
\end{lstlisting}

\begin{tcbinfo}{เคล็ดลับ: เมื่อไหร่ใช้แบบไหน?}
กฎจากประสบการณ์: ขั้นตอนหลักใช้ NavigationStack ดันหน้า (รายการ รายละเอียด), การทำงานชั่วคราวใช้ป๊อปอัป sheet (เพิ่ม, แก้ไข, ยืนยัน), ส่วนบล็อกใหญ่ที่อยู่เคียงกันใช้ TabView (แท็บด้านล่าง) ประกอบสามอย่างนี้เข้าด้วยกัน ก็คือโครงของแอปเกือบทุกตัว

\end{tcbinfo}

\booksection{List: 'ลิสต์แบบฉบับ iOS'}

คอมโพเนนต์ที่ใช้บ่อยที่สุดและมี'กลิ่นอาย Apple'ที่สุดบน iOS ก็คือ List — เส้นแบ่งรายการ ปัดเพื่อลบ ไฮไลต์เมื่อแตะ มาให้ตั้งแต่แกะกล่อง:

\begin{lstlisting}[language=Swift]
struct RecordListView: View {
    @State private var records = ["มื้อกลางวัน 25.5", "รถไฟใต้ดิน 6.0", "เงินเดือน +8000"]

    var body: some View {
        List {
            ForEach(records, id: \.self) { r in
                Text(r)
            }
            .onDelete { indexSet in          // ปัดซ้ายเพื่อลบ
                records.remove(atOffsets: indexSet)
            }
        }
    }
}
\end{lstlisting}

รายการ + NavigationLink + sheet + TabView — เอาสี่อย่างนี้มาประกอบกัน คุณก็สร้างโครงแอปที่สมบูรณ์แบบ'ลิสต์สมุดบัญชี แตะดูรายละเอียด กด + แล้วฟอร์มบันทึกเด้งขึ้นมา'ได้แล้ว โปรเจกต์สมุดบัญชีในบทหน้าก็สร้างด้วยวิธีนี้แหละ

\begin{tcbpractice}{ลงมือทำ: สร้าง'โครงสามหน้า'}
\booksection{วงจรชีวิตของแอป: ฟอร์กราวด์กับแบ็กกราวด์}

คล้ายวงจรชีวิต Activity ของ Android แอป iOS ก็มีการเปลี่ยนสถานะแบบ'เกิด, อยู่ฟอร์กราวด์, ลงแบ็กกราวด์, ตาย'เช่นกัน ใน SwiftUI ให้เฝ้าดูผ่านค่าสภาพแวดล้อม scenePhase:

\begin{lstlisting}[language=Swift]
@main
struct MyApp: App {
    @Environment(\.scenePhase) private var phase

    var body: some Scene {
        WindowGroup {
            ContentView()
        }
        .onChange(of: phase) { _, newPhase in
            switch newPhase {
            case .active:    print("กลับมาพื้นหน้า กู้การรีเฟรชต่อ")
            case .inactive:  print("อยู่ระหว่างเปลี่ยน (สายเข้า/พึ่งสลับไปพื้นหลัง)")
            case .background: print("ลงพื้นหลัง บันทึกฉบับร่าง")
            @unknown default: break
            }
        }
    }
}
\end{lstlisting}

การใช้งานทั่วไป: กลับมาที่ฟอร์กราวด์ก็รีเฟรชข้อมูล ลงแบ็กกราวด์ก็บันทึกฉบับร่าง iOS ควบคุมการรันพื้นหลังเข้มงวดมาก — แอปทั่วไปพอเข้าแบ็กกราวด์จะถูกระงับ (แช่แข็ง) ในเวลาไม่นาน อย่าหวังให้แอปแอบทำงานหนักเบื้องหลัง ถ้าจำเป็นต้องทำงานตามเวลาจริง ๆ ให้ใช้การแจ้งเตือนในเครื่องจากบทที่ 68 ไม่ใช่การ'อาศัยอยู่เบื้องหลังถาวร'

สร้างโครงสมุดบัญชีเวอร์ชันข้อมูลปลอม: TabView สามแท็บ (สมุดบัญชี/สถิติ/ตั้งค่า) หน้าสมุดบัญชีใช้ List แสดงข้อมูลปลอม แตะแถวแล้วเปลี่ยนไปหน้ารายละเอียด (NavigationStack) ปุ่ม + มุมขวาบนเด้ง sheet บันทึกรายการขึ้นมา ไม่เก็บข้อมูล ไม่ต่อเน็ต เอาไว้ฝึกสร้างโครงล้วน ๆ — ทำโครงนี้ได้ พอถึงบทที่ 65 ใส่ข้อมูลจริงลงไปก็ลื่นไหลเป็นธรรมชาติ

\end{tcbpractice}

\begin{tcbsummary}{สรุปบทนี้}
\begin{itemize}
  \item เครื่องมือสเตตสี่ตัว: @State (ส่วนตัวของหน้า), @Binding (ยืมมาใช้), @Observable (โลจิก), .environment (ทั้งแอป)
  \item โปรเจกต์ใหม่ใช้มาโคร @Observable เขียนวิวโมเดล; @ObservableObject ในบทเรียนเก่า อ่านออกก็พอ
  \item NavigationStack ดันหน้าใหม่; @Environment(\textbackslash{}.dismiss) ใช้ย้อนกลับด้วยโค้ด
  \item ขั้นตอนหลักดันหน้า งานชั่วคราวเด้ง sheet บล็อกใหญ่เคียงกันใช้ TabView
  \item List คือ'ลิสต์แบบฉบับ iOS': เส้นแบ่งและปัดเพื่อลบ ใช้ได้ทันทีแบบแกะกล่อง
\end{itemize}

\end{tcbsummary}
